% Actual class: 04
% Slide decks: M1-L7 LAN-Ethernet-and-Its-Evolutions (2026); M1-L8 LAN-WLAN (2026)
% Exact scope: L7 pp. 1--43; L8 pp. 1--47
% NTULearn supplement: classroom recording transcript, used only to verify coverage and emphasis
% Human verified: Yes

\section{第 04 节课：Ethernet 与 WLAN}
\label{lecture:SC2008-L04}

\lecturemeta
  {SC2008-L04}
  {2026-09-01}
  {M1-L7 LAN-Ethernet-and-Its-Evolutions (2026)；M1-L8 LAN-WLAN (2026)}
  {L7 pp.~1--43；L8 pp.~1--47}
  {课堂录像字幕（用于确认两讲完整覆盖及 examinable emphasis；原文未入库）}
  {已确认（2026-09-01）}

\subsection{本讲主线与考试权重}

Ethernet（以太网）与 IEEE 802.11 WLAN（无线局域网）都解决 shared medium（共享介质）的接入问题，但物理条件决定了不同策略：经典 shared Ethernet 能在发送时检测 collision，采用 CSMA/CD 与 BEB；WLAN 受 attenuation（衰减）和 hidden terminals（隐藏终端）影响，无法可靠检测 collision，改用 CSMA/CA 与可选的 RTS/CTS reservation。MARP 再把 contention cost（竞争开销）摊销到较长的 collision-free data phase。

课堂录像明确将 Ethernet 的 MAC performance calculation 与 WLAN 的 MARP calculation 作为主要 examinable 内容；history、physical variants、802.11 physical layer 与 protocol evolution 只保留理解核心机制所需的结论。

\lecturetopic{SC2008-L04-T01}{Ethernet Frame 与 MAC Addressing}

\begin{figure}[htbp]
  \centering
  \resizebox{0.98\textwidth}{!}{\begin{tikzpicture}[x=0.82cm,y=0.9cm,>=Latex,font=\small]
  \def\h{1.1}
  \draw[fill=gray!12]    (0,0) rectangle (2.2,\h) node[midway] {Preamble};
  \draw[fill=gray!18]    (2.2,0) rectangle (3.2,\h) node[midway] {SFD};
  \draw[fill=blue!12]    (3.2,0) rectangle (4.8,\h) node[midway] {DA};
  \draw[fill=blue!12]    (4.8,0) rectangle (6.4,\h) node[midway] {SA};
  \draw[fill=orange!18]  (6.4,0) rectangle (7.8,\h) node[midway] {L/T};
  \draw[fill=green!13]   (7.8,0) rectangle (12.2,\h) node[midway] {Data + Padding};
  \draw[fill=red!12]     (12.2,0) rectangle (13.7,\h) node[midway] {FCS};

  \foreach \x/\lab in {1.1/7 B,2.7/1 B,4.0/6 B,5.6/6 B,7.1/2 B,10.0/46--1500 B,12.95/4 B}
    \node[below,font=\scriptsize] at (\x,-0.08) {\lab};

  \draw[<->,thick,noteBlue] (3.2,1.48) -- (13.7,1.48)
    node[midway,above,font=\scriptsize] {MAC frame: 64--1518 B};
  \draw[<->,thick,black!65] (0,-0.65) -- (3.2,-0.65)
    node[midway,below,font=\scriptsize] {synchronization; excluded from MAC-frame length};
\end{tikzpicture}
}
  \caption{Classic Ethernet frame format。课程的 64--1518 B MAC-frame length 不包括 Preamble 与 SFD。}
  \label{fig:sc2008-l04-ethernet-frame}
\end{figure}

\begin{center}
\small
\begin{tabularx}{\textwidth}{@{}>{\raggedright\arraybackslash}p{2.9cm}>{\raggedright\arraybackslash}p{2.2cm}X@{}}
  \toprule
  Field & Length & Function \\
  \midrule
  Preamble + SFD & $7+1$ B & Receiver clock synchronization 与 start-frame delimiter；不计入 MAC-frame length \\
  DA + SA & $6+6$ B & 48-bit destination/source MAC addresses \\
  Length/Type & 2 B & 指示 payload length 或 upper-layer protocol，完成 demultiplexing \\
  Data + Padding & 46--1500 B & Payload；不足 46 B 时 padding 以满足 minimum frame size \\
  FCS & 4 B & 对 DA 至 Data/Padding 执行 CRC error detection \\
  \bottomrule
\end{tabularx}
\end{center}

因此 classic untagged Ethernet 有
\[
  \boxed{L_{\min}=64\ \text{B}},
  \qquad
  \boxed{L_{\max}=1518\ \text{B}},
  \qquad
  \boxed{\text{payload}_{\max}=1500\ \text{B}}.
\]
不要把 maximum payload 误写成 whole-frame maximum。

Destination address 的 I/G bit 是 first octet 的 least significant bit：$0$ 表示 individual/unicast，$1$ 表示 group/multicast；\texttt{FF:FF:FF:FF:FF:FF} 是 broadcast。书写成 hexadecimal 时，可检查 first octet 的第二个 hex digit：偶数为 unicast，奇数为 multicast，但应先排除全 \texttt{FF} 的 broadcast。

\textbf{对应例题：}\relatedexercise{SC2008-L04-E01}{讲义例题：MAC Address Classification}。

\lecturetopic{SC2008-L04-T02}{CSMA/CD、64-byte Minimum 与 Jam Signal}

Classic shared half-duplex Ethernet 使用 1-persistent CSMA/CD：channel idle 时立即发送；busy 时持续监听并在变 idle 后立即发送；transmission 期间继续检测 collision。检测到 collision 后 abort frame、发送 jam signal，再进入 random backoff。现代 switched full-duplex Ethernet 的 links 不共享 collision domain，通常不再运行 CSMA/CD，但 frame format 与 64-byte minimum 仍被保留。

若 maximum one-way propagation delay 为 $\tau$，data rate 为 $R$，sender 必须在最坏 collision indication 返回前仍在发送：
\[
  T_{\mathrm{frame}}=\frac{L}{R}\ge 2\tau,
  \qquad
  \boxed{L_{\min}\ge2R\tau}.
\]

\begin{figure}[htbp]
  \centering
  \resizebox{0.86\textwidth}{!}{\begin{tikzpicture}[>=Latex,every node/.style={font=\small}]
  \coordinate (A0) at (0,0);
  \coordinate (B0) at (9,0);
  \coordinate (Atwo) at (0,-5.2);
  \coordinate (Btwo) at (9,-5.2);

  \draw[->,thick] (A0) -- (Atwo) node[below] {time};
  \draw[->,thick] (B0) -- (Btwo) node[below] {time};
  \node[above] at (A0) {Station A};
  \node[above] at (B0) {Station B};
  \draw[black!55] (A0) -- (B0) node[midway,above,font=\scriptsize] {maximum separation};

  \draw[->,very thick,noteBlue] (0,-0.35) -- (8.75,-2.55)
    node[midway,above,sloped] {A's signal};
  \fill[ntuRed] (8.75,-2.55) circle[radius=2pt];
  \node[anchor=east,text=ntuRed,font=\scriptsize] at (8.60,-3.25) {collision near B};
  \draw[->,very thick,ntuRed] (8.75,-2.55) -- (0,-4.75)
    node[midway,above,sloped] {collision indication};
  \node[anchor=west,align=left] at (9.25,-1.75)
    {B starts just before\\A's signal arrives};
  \draw[->,thick,black!75] (9.25,-2.15) -- (8.82,-2.50);

  \draw[dashed] (-0.15,-2.55) -- (9.15,-2.55);
  \node[anchor=east] at (-0.2,-2.55) {$\tau$};
  \draw[dashed] (-0.15,-4.75) -- (9.15,-4.75);
  \node[anchor=east] at (-0.2,-4.75) {$2\tau$};
  \node[anchor=west,text=noteBlue] at (0.15,-0.35) {A begins transmission};
  \node[anchor=west,text=ntuRed,fill=white,inner sep=1.5pt] at (0.30,-4.15)
    {A detects collision};
\end{tikzpicture}
}
  \caption{Worst-case collision detection：远端在源信号到达前一瞬间发送，源端最迟在 $2\tau$ 检测 collision。}
  \label{fig:sc2008-l04-csma-cd-timing}
\end{figure}

Classic 10 Mbps Ethernet 规定 slot time 为 $51.2\ \mu\mathrm{s}$，因此
\[
  L_{\min}=(10\times10^6)(51.2\times10^{-6})
  =512\ \text{bits}=64\ \text{B}.
\]
其中 64 B 可写成 $14$ B header $+46$ B minimum data/pad $+4$ B FCS。Jam signal 用来使 collision 足够明显，而不是修复损坏 frame；本讲计算按 48-bit jam signal 处理。

\textbf{已学理论例题：}\relatedexercise{SC2008-L03-E03}{CSMA/CD Minimum Frame Length}。

\lecturetopic{SC2008-L04-T03}{Binary Exponential Backoff 与 Ethernet Domains}

第 $n$ 次 collision 后，BEB（Binary Exponential Backoff，二进制指数退避）令 station 独立均匀选择
\[
  K\in\{0,1,\ldots,2^n-1\},
  \qquad
  \boxed{T_{\mathrm{backoff}}=K\times51.2\ \mu\mathrm{s}}.
\]
Exponent 在 $n=10$ 后封顶，故最大 $K=1023$；一个 frame 通常在 16 次失败后放弃。$K$ 决定当前 retrial 的随机等待，16 次是 maximum retry count，两者不能混淆。对两个 stations、$m$ 个等概率 slots，二者再次选中相同 slot 的概率为
\[
  \boxed{P(\text{collision})=\frac{1}{m}}.
\]
BEB 先用较小 range 追求低 delay；若 collision 持续则扩大 range，逐渐以更多 idle time 换取更高 collision-resolution probability。

例如 third collision 后 $K\in\{0,\ldots,7\}$：指定 station 选到 $K=7$ 的概率为 $1/8$，maximum backoff delay 为 $7(51.2\ \mu\mathrm{s})=358.4\ \mu\mathrm{s}$。这与 ``最多尝试 16 次'' 是两个不同概念。

\textbf{对应例题：}\relatedexercise{SC2008-L04-E02}{讲义例题：BEB Retrial Probability}。

\begin{center}
\small
\begin{tabularx}{\textwidth}{@{}>{\raggedright\arraybackslash}p{3.0cm}XX@{}}
  \toprule
  Device & Collision domain & Broadcast domain \\
  \midrule
  Repeater / Hub & 不隔离；向其他 ports 重复 signal & 不隔离 \\
  Bridge / Layer-2 Switch & 每个 segment/port 隔离 collision；full-duplex link 可完全避免 collision & 默认仍转发 broadcast；同一 VLAN 通常仍属一个 broadcast domain \\
  Router & 隔离 & 通常隔离；是否转发特殊 broadcast 取决于 configuration \\
  \bottomrule
\end{tabularx}
\end{center}

\lecturetopic{SC2008-L04-T04}{WLAN、CSMA/CA 与 Hidden/Exposed Terminals}

Infrastructure mode（基础设施模式）让 stations 经 AP（Access Point，接入点）通信；一个 AP 与覆盖区内 stations 构成 BSS（Basic Service Set），多个 BSS 经 distribution system 连接成 ESS（Extended Service Set）。Ad hoc mode（自组织模式）则允许 stations 直接组成 distributed network。

Wireless link 会受到 path loss、interference 与 multipath propagation。Transmitter 的本地发射信号远强于远端 received signal，加上 hidden terminal，使 station 无法在 transmission 期间可靠判断 receiver 处是否 collision，因此 IEEE 802.11 DCF 使用 CSMA/CA 而非 CSMA/CD。

\begin{figure}[htbp]
  \centering
  \resizebox{0.96\textwidth}{!}{\begin{tikzpicture}[x=1cm,y=1cm,>=Latex,font=\small]
  \tikzset{
    station/.style={draw,rounded corners,minimum width=1.05cm,minimum height=0.68cm,fill=blue!10},
    ap/.style={draw,very thick,circle,minimum size=0.82cm,fill=orange!18}
  }

  \node[station] (A) at (0,3.9) {A};
  \node[ap]      (B) at (5.3,3.9) {B};
  \node[station] (C) at (10.6,3.9) {C};
  \draw[<->,thick,noteBlue] (A) -- (B) node[midway,above] {hear each other};
  \draw[<->,thick,noteBlue] (B) -- (C) node[midway,above] {hear each other};
  \draw[dashed,thick,ntuRed] (A.south east) to[bend right=24]
    node[midway,below,text=ntuRed] {A and C cannot hear each other} (C.south west);

  \coordinate (At) at (0,1.65);
  \coordinate (Bt) at (5.3,1.65);
  \coordinate (Ct) at (10.6,1.65);
  \draw[->] (At) -- +(0,-4.2) node[below] {time};
  \draw[->] (Bt) -- +(0,-4.2) node[below] {time};
  \draw[->] (Ct) -- +(0,-4.2) node[below] {time};
  \node[above] at (At) {A};
  \node[above] at (Bt) {B};
  \node[above] at (Ct) {C};

  \draw[->,very thick,noteBlue] (0,1.25) -- (5.3,0.75) node[midway,above,sloped] {RTS};
  \draw[->,very thick,orange!85!black] (5.3,0.35) -- (0,-0.15) node[midway,above,sloped] {CTS};
  \draw[->,very thick,orange!85!black] (5.3,0.35) -- (10.6,-0.15) node[midway,above,sloped] {CTS};
  \draw[very thick,gray!55] (10.6,-0.15) -- (10.6,-2.85);
  \node[anchor=west,align=left,font=\scriptsize] at (10.8,-1.45) {NAV: C defers};
  \draw[->,very thick,green!55!black] (0,-0.70) -- (5.3,-1.30) node[midway,above,sloped] {DATA};
  \draw[->,very thick,ntuRed] (5.3,-1.85) -- (0,-2.35) node[midway,above,sloped] {ACK};
\end{tikzpicture}
}
  \caption{Hidden terminal 与 RTS--CTS--DATA--ACK：B 的 CTS 使听得到 B 的其他 stations 设置 NAV 并 defer。}
  \label{fig:sc2008-l04-hidden-terminal}
\end{figure}

若 A、C 互相听不到但都能向 B 发送，两者可能同时判断 channel idle，并在 B 处 collision。RTS/CTS 先用 short control frames 预约：sender 发 RTS，receiver 广播 CTS；其他听到 CTS 的 stations 按 NAV（Network Allocation Vector）保持 silent，随后完成 DATA 与 ACK。RTS 仍可能 collision，因此 collision avoidance 不是 collision elimination，只是尽量避免 long data-frame collision。

Exposed terminal（暴露终端）是 opportunity wasting：C 因听到 A$\to$B 而停止向 D 发送，即使 C$\to$D 与 A$\to$B 本可并行且不会在各自 receiver 处 collision。RTS/CTS 能缓解 hidden terminal，却不能完全恢复这类 spatial reuse。

802.11 frame 可使用最多四个 address fields；ordinary infrastructure forwarding 常使用三个地址，把 wireless receiver/transmitter 与 distribution-system destination 分开。该格式和 physical generations 在本讲标为 non-examinable。

\lecturetopic{SC2008-L04-T05}{MARP Utilization 与 Amortization}

MARP（Multi-Access Reservation Protocol，多路访问预约协议）把 transmission window 分为 contention-based reservation 与 collision-free data transmission 两个 phases。

\begin{figure}[htbp]
  \centering
  \resizebox{0.94\textwidth}{!}{\begin{tikzpicture}[x=1.15cm,y=1cm,>=Latex,font=\small]
  \node[anchor=west,font=\bfseries] at (0,2.1) {Phase 1: Reservation};
  \node[anchor=west,font=\bfseries] at (6.0,2.1) {Phase 2: Data Transmission};

  \foreach \i in {0,1,2}{
    \draw[fill=red!10] (\i*1.35,0.55) rectangle +(1.05,0.85);
  }
  \node at (0.525,0.975) {fail};
  \node at (1.875,0.975) {fail};
  \node at (3.225,0.975) {fail};
  \draw[fill=green!18] (4.05,0.55) rectangle +(1.05,0.85) node[midway] {success};
  \draw[fill=blue!12] (6.0,0.55) rectangle (12.3,1.4) node[midway] {collision-free data};

  \draw[<->] (0,0.12) -- (1.05,0.12) node[midway,below] {$v$};
  \draw[<->] (0,-0.72) -- (5.10,-0.72)
    node[midway,below] {$E[X]=1/S_r$ trials, expected time $v/S_r$};
  \draw[<->] (6.0,0.12) -- (12.3,0.12) node[midway,below] {$u$};
  \node[align=center,font=\scriptsize] at (6.15,-1.55)
    {$S=\dfrac{u}{u+v/S_r}$ when reservation carries no message data};
\end{tikzpicture}
}
  \caption{MARP transmission window：平均需 $1/S_r$ 次 reservation trials，reservation overhead 为 $v/S_r$。}
  \label{fig:sc2008-l04-marp-window}
\end{figure}

设每个 reservation trial 成功概率为 $S_r$、长度为 $v$，直到第一次成功所需次数为 $X$。在 independent identical trials（独立同分布试验）假设下，
\[
  P(X=k)=S_r(1-S_r)^{k-1},
  \qquad
  \boxed{E[X]=\frac{1}{S_r}},
  \qquad
  \boxed{E[T_{\mathrm{reservation}}]=\frac{v}{S_r}}.
\]

若 message transmission time 为 $u$：
\[
\boxed{
  S=
  \begin{cases}
    \displaystyle\frac{u}{u+v/S_r},
      & \text{reservation 不携带 message data},\\[0.9em]
    \displaystyle\frac{u}{(u-v)+v/S_r},
      & \text{成功的 reservation 携带 $v$ 单位 message data}.
  \end{cases}}
\]
第二式假设 $u\ge v$，失败 trials 不贡献 useful data，只有成功 reservation 中的 $v$ 计入 message。增大 $u$ 或 $S_r$、减小 $v$ 都可提高 utilization；本质是把短 reservation cost amortize（摊销）到较长的 collision-free phase。

\textbf{对应例题：}\relatedexercise{SC2008-L04-E03}{讲义例题：MARP Utilization}。

\subsection{一分钟回顾}

\[
\boxed{
\begin{gathered}
L_{\mathrm{Ethernet}}=64\text{--}1518\ \text{B},\qquad
L_{\min}\ge2R\tau,\\
T_{\mathrm{backoff}}=K(51.2\ \mu\mathrm{s}),\quad
K\in\{0,\ldots,2^n-1\},\\
E[X_{\mathrm{MARP}}]=1/S_r,\qquad
S_{\mathrm{MARP}}=\frac{u}{u+v/S_r}
\quad\text{(no data in reservation)}.
\end{gathered}}
\]
